\documentclass[aspectratio=169]{beamer}

% ─── Tema y color UNS ────────────────────────────────────────────────────────
\usetheme{Madrid}
\definecolor{unsblue}{RGB}{30,60,120}
\definecolor{unslightblue}{RGB}{220,230,245}
\definecolor{alertred}{RGB}{180,30,30}
\usecolortheme[named=unsblue]{structure}
\setbeamercolor{palette primary}{bg=unsblue, fg=white}
\setbeamercolor{palette secondary}{bg=unsblue!80, fg=white}
\setbeamercolor{palette tertiary}{bg=unsblue!60, fg=white}
\setbeamercolor{palette quaternary}{bg=unsblue, fg=white}
\setbeamercolor{block title}{bg=unsblue, fg=white}
\setbeamercolor{block body}{bg=unslightblue}
\setbeamercolor{block title alerted}{bg=alertred, fg=white}
\setbeamercolor{block body alerted}{bg=red!8}
\setbeamercolor{block title example}{bg=unsblue!60!black, fg=white}
\setbeamercolor{block body example}{bg=unsblue!8}

% ─── Paquetes ────────────────────────────────────────────────────────────────
\usepackage[utf8]{inputenc}
\usepackage[T1]{fontenc}
\usepackage{babel}
\babelprovide[main, import]{spanish}
\usepackage{amsmath, amssymb}
\usepackage{bm}
\usepackage{booktabs}
\usepackage{graphicx}

% ─── Comandos ────────────────────────────────────────────────────────────────
\newcommand{\R}{\mathbb{R}}
\newcommand{\bx}{\bm{x}}
\newcommand{\E}{\mathbb{E}}

% ─── Metadatos ───────────────────────────────────────────────────────────────
\title[MAD1 -- Unidad 4]{Aprendizaje Supervisado}
\subtitle{Clase 20: Introducción y Árboles de Decisión}
\author[J. Bavio]{Dr.\ José Bavio}
\institute[UNS]{Departamento de Matemática\\Universidad Nacional del Sur}
\date{2026}

% ─────────────────────────────────────────────────────────────────────────────
\begin{document}
% ─────────────────────────────────────────────────────────────────────────────

\begin{frame}
  \titlepage
\end{frame}

\begin{frame}{Contenido de la clase}
  \tableofcontents
\end{frame}

% ═════════════════════════════════════════════════════════════════════════════
\section{¿Qué es el Aprendizaje Supervisado?}
% ═════════════════════════════════════════════════════════════════════════════

\begin{frame}{¿Qué es el Aprendizaje Supervisado?}
  \begin{block}{Idea central}
    Tenemos datos \textbf{etiquetados}: para cada observación $\bx_i$ conocemos la respuesta $y_i$.\\
    El objetivo es aprender una función $f$ tal que $f(\bx) \approx y$ para datos nuevos.
  \end{block}

  \medskip
  \begin{exampleblock}{Dicho de otra forma}
    Es como aprender con un \textbf{solucionario}: el modelo estudia ejemplos resueltos
    y luego aplica lo aprendido a ejercicios que nunca vio.
  \end{exampleblock}

  \medskip
  Según el tipo de $y$:
  \begin{itemize}
    \item $y \in \R$ \quad$\Rightarrow$\quad \textbf{regresión} (Unidad 3)
    \item $y \in \{1, \ldots, K\}$ \quad$\Rightarrow$\quad \textbf{clasificación} (esta unidad)
  \end{itemize}
\end{frame}

\begin{frame}{El objetivo real: generalización}
  \begin{alertblock}{Error de entrenamiento $\neq$ error real}
    Memorizar los datos de entrenamiento no sirve.\\
    Lo que importa es predecir bien datos \textbf{que el modelo nunca vio}.
  \end{alertblock}

  \medskip
  \begin{exampleblock}{Analogía}
    Un alumno que memoriza las respuestas del parcial anterior sin entender
    los conceptos va a reprobar el parcial nuevo.\\
    El modelo que memoriza los datos de entrenamiento hace exactamente eso.
  \end{exampleblock}

  \medskip
  \textbf{Error de generalización:}
  \[
    \text{Err} = \E\bigl[\ell(y,\, f(\bx))\bigr]
  \]
  donde $\ell$ es la función de pérdida (ej.: pérdida 0-1 para clasificación).
\end{frame}

% ─────────────────────────────────────────────────────────────────────────────
\section{Sesgo, Varianza y Overfitting}
% ─────────────────────────────────────────────────────────────────────────────

\begin{frame}{Sesgo y Varianza}
  Para regresión, el error de predicción descompone en:
  \[
    \text{ECM} = \underbrace{\text{Var}[\hat{f}(\bx_0)]}_{\text{varianza}}
              + \underbrace{[\text{Sesgo}(\hat{f}(\bx_0))]^2}_{\text{sesgo}^2}
              + \underbrace{\sigma^2_\varepsilon}_{\text{irreducible}}
  \]

  \medskip
  \begin{columns}
    \column{0.48\textwidth}
    \begin{block}{Alto sesgo (underfitting)}
      El modelo es \textbf{demasiado simple}.\\
      No captura la estructura real.\\[4pt]
      \textit{Como ajustar una recta a datos claramente curvos.}
    \end{block}

    \column{0.48\textwidth}
    \begin{block}{Alta varianza (overfitting)}
      El modelo es \textbf{demasiado complejo}.\\
      Aprende el ruido de los datos.\\[4pt]
      \textit{Como un modelo que conecta todos los puntos con una curva serpenteante.}
    \end{block}
  \end{columns}
\end{frame}

\begin{frame}{El compromiso sesgo-varianza}
  \begin{center}
    \begin{tabular}{ccc}
      \toprule
      \textbf{Complejidad} & \textbf{Sesgo} & \textbf{Varianza} \\
      \midrule
      Baja  & Alto  & Baja \\
      Media & Medio & Media \\
      Alta  & Bajo  & Alta  \\
      \bottomrule
    \end{tabular}
  \end{center}

  \medskip
  \begin{exampleblock}{En la práctica}
    El error de \textbf{validación} tiene forma de U en función de la complejidad.
    El mínimo de esa curva es el punto óptimo.
    El error de \textbf{entrenamiento} baja siempre: no sirve como indicador.
  \end{exampleblock}
\end{frame}

% ─────────────────────────────────────────────────────────────────────────────
\section{Validación Cruzada}
% ─────────────────────────────────────────────────────────────────────────────

\begin{frame}{Partición de los datos}
  \begin{block}{Tres conjuntos}
    \begin{itemize}
      \item \textbf{Entrenamiento:} ajuste de los parámetros del modelo.
      \item \textbf{Validación:} selección de hiperparámetros y comparación de modelos.
      \item \textbf{Prueba (test):} evaluación final. Se toca \textbf{una única vez}.
    \end{itemize}
  \end{block}

  \medskip
  \begin{alertblock}{Regla de oro}
    El conjunto de prueba no debe influir en ninguna decisión de modelado.\\
    Mirarlo varias veces equivale a ``hacer trampa'' con ese conjunto.
  \end{alertblock}
\end{frame}

\begin{frame}{Validación cruzada $k$-fold}
  \textbf{Idea:} dividir el conjunto de entrenamiento en $k$ partes iguales (\textit{folds}).

  \medskip
  \begin{enumerate}
    \item Dividir en $k$ folds.
    \item Para $j = 1, \ldots, k$: entrenar con los otros $k-1$ folds, evaluar en el fold $j$.
    \item Promediar los $k$ errores:
    \[
      \text{CV}_k = \frac{1}{k}\sum_{j=1}^{k} \text{Err}_j
    \]
  \end{enumerate}

  \medskip
  \begin{exampleblock}{Coloquialmente}
    Es como rendir el mismo examen 5 veces con distintos grupos de alumnos para
    tener una estimación más confiable de cuánto aprendió el modelo.
  \end{exampleblock}

  \medskip
  \textbf{Valores típicos:} $k = 5$ o $k = 10$.\\
  \textbf{Caso extremo:} $k = n$ (\textit{leave-one-out}, LOOCV).
\end{frame}

\begin{frame}{CV estratificada}
  \begin{block}{Problema con clases desbalanceadas}
    Si el 95\% de los datos son clase A y solo el 5\% clase B,
    un fold podría no tener ningún ejemplo de clase B.
  \end{block}

  \medskip
  \begin{exampleblock}{Solución: CV estratificada}
    Cada fold mantiene la misma \textbf{proporción de clases} que el conjunto original.\\[4pt]
    Es lo que hace \texttt{StratifiedKFold} en scikit-learn.
    Siempre usarla en clasificación.
  \end{exampleblock}
\end{frame}

% ─────────────────────────────────────────────────────────────────────────────
\section{Árboles de Decisión}
% ─────────────────────────────────────────────────────────────────────────────

\begin{frame}{¿Qué es un árbol de decisión?}
  \begin{block}{Idea}
    Particionar el espacio de features con reglas del tipo\\
    \centerline{``$x_j \leq t$ \, ¿sí o no?''}
    hasta llegar a regiones lo más \textbf{puras} posible.
  \end{block}

  \medskip
  \begin{exampleblock}{Analogía cotidiana}
    Es exactamente el juego del \textbf{Akinator} o el ``animal, vegetal o mineral'':
    preguntas binarias sucesivas que van acotando las posibilidades
    hasta llegar a una respuesta.
  \end{exampleblock}

  \medskip
  Estructura:
  \begin{itemize}
    \item \textbf{Nodo raíz:} todos los datos.
    \item \textbf{Nodos internos:} cada uno hace una pregunta $x_j \leq t$.
    \item \textbf{Hojas:} regiones terminales que asignan una clase.
  \end{itemize}
\end{frame}

\begin{frame}{Criterios de partición: Gini}
  \textbf{Índice de Gini} para un nodo con proporciones de clase $p_1, \ldots, p_K$:
  \[
    G = 1 - \sum_{k=1}^{K} p_k^2
  \]

  \begin{itemize}
    \item $G = 0$: nodo \textbf{puro} (todas las observaciones son de la misma clase). Ideal.
    \item $G$ máximo: clases perfectamente mezcladas. Nodo inútil.
  \end{itemize}

  \medskip
  \begin{exampleblock}{Intuición}
    Gini mide la probabilidad de equivocarse si clasificamos al azar según las proporciones del nodo.
    Un nodo puro tiene probabilidad de error 0, o sea, $G = 0$.
  \end{exampleblock}
\end{frame}

\begin{frame}{Criterios de partición: Entropía}
  \textbf{Entropía} (ganancia de información):
  \[
    H = -\sum_{k=1}^{K} p_k \log_2(p_k)
  \]

  La \textbf{ganancia de información} de un split es la caída de entropía:
  \[
    \Delta H = H(\text{padre}) - \frac{n_L}{n} H(\text{izq.}) - \frac{n_R}{n} H(\text{der.})
  \]

  \medskip
  \begin{block}{¿Gini o entropía?}
    En la práctica producen resultados casi idénticos.
    Gini es más rápido (no requiere logaritmos).
    scikit-learn usa Gini por defecto.
  \end{block}
\end{frame}

\begin{frame}{Algoritmo CART}
  \textbf{CART} (\textit{Classification and Regression Trees}) construye el árbol de manera \textit{greedy}:

  \medskip
  \begin{enumerate}
    \item Partir de todos los datos en el nodo raíz.
    \item Para cada feature $j$ y cada umbral $t$: calcular la impureza del split.
    \item Elegir el par $(j^*, t^*)$ que minimiza la impureza total.
    \item Repetir en cada subárbol (recursión).
    \item Parar cuando se cumple un criterio de parada.
  \end{enumerate}

  \medskip
  \begin{alertblock}{Greedy $\neq$ óptimo global}
    El árbol construido por CART no es necesariamente el mejor árbol posible.
    Elegir el mejor split en cada nodo no garantiza el mejor árbol completo.
    Es como elegir siempre la jugada que parece mejor en ajedrez
    sin pensar en las próximas jugadas.
  \end{alertblock}
\end{frame}

\begin{frame}{Criterios de parada e hiperparámetros}
  \begin{block}{Cuándo dejar de crecer el árbol}
    \begin{itemize}
      \item \texttt{max\_depth}: profundidad máxima del árbol.
      \item \texttt{min\_samples\_split}: mínimo de muestras para dividir un nodo.
      \item \texttt{min\_samples\_leaf}: mínimo de muestras en una hoja.
      \item \texttt{max\_leaf\_nodes}: máximo de hojas totales.
    \end{itemize}
  \end{block}

  \medskip
  \begin{exampleblock}{Intuición}
    Sin restricciones, el árbol crece hasta que cada hoja tiene un solo dato:
    accuracy 100\% en entrenamiento, pero generaliza pésimo.
    Los criterios de parada son la ``poda preventiva''.
  \end{exampleblock}
\end{frame}

\begin{frame}{Poda por complejidad de costo}
  Otra estrategia: dejar crecer el árbol al máximo y luego \textbf{podarlo}.

  \medskip
  Se busca el subárbol $T$ que minimiza:
  \[
    C_\alpha(T) = \underbrace{\text{error de clasificación}}_{\text{en entrenamiento}} + \alpha \cdot \underbrace{|T|}_{\text{n° de hojas}}
  \]

  \begin{itemize}
    \item $\alpha = 0$: árbol completo, sin penalización.
    \item $\alpha$ grande: árbol muy pequeño (solo la raíz si $\alpha \to \infty$).
    \item El $\alpha$ óptimo se elige por validación cruzada.
  \end{itemize}

  \medskip
  \begin{exampleblock}{Analogía}
    Es como la regulación $L_1/L_2$ en regresión, pero para árboles:
    penalizamos la complejidad del modelo para evitar el sobreajuste.
  \end{exampleblock}
\end{frame}

\begin{frame}{Ventajas y limitaciones de los árboles}
  \begin{columns}
    \column{0.48\textwidth}
    \begin{block}{Ventajas}
      \begin{itemize}
        \item Muy interpretables:\\``Si $x_1 > 3$ y $x_2 \leq 7$, entonces clase A''
        \item No requieren normalización
        \item Manejan variables mixtas
        \item Capturan interacciones
      \end{itemize}
    \end{block}

    \column{0.48\textwidth}
    \begin{alertblock}{Limitaciones}
      \begin{itemize}
        \item Alta varianza: cambios pequeños en los datos generan árboles muy distintos
        \item Fronteras de decisión solo axiales (paralelas a los ejes)
        \item Tendencia al sobreajuste
      \end{itemize}
    \end{alertblock}
  \end{columns}

  \medskip
  \begin{exampleblock}{Por eso se usan como base de ensambles}
    Un árbol solo es inestable. Muchos árboles juntos son robustos.\\
    Eso es exactamente Random Forest (próxima clase).
  \end{exampleblock}
\end{frame}

% ─────────────────────────────────────────────────────────────────────────────
\begin{frame}{Resumen de la clase}
  \begin{itemize}
    \item El \textbf{aprendizaje supervisado} aprende de ejemplos etiquetados para predecir datos nuevos.
    \item El objetivo es la \textbf{generalización}, no el ajuste perfecto al entrenamiento.
    \item El compromiso \textbf{sesgo-varianza} explica el underfitting y el overfitting.
    \item La \textbf{validación cruzada $k$-fold} estima el error de generalización.
    \item Los \textbf{árboles de decisión} particionan el espacio con reglas binarias usando Gini o entropía.
    \item El algoritmo \textbf{CART} es greedy; la poda controla la complejidad.
  \end{itemize}

  \bigskip
  \textbf{Próxima clase:} Random Forest --- cómo convertir árboles inestables en un modelo robusto.
\end{frame}

\end{document}
